% Generated draft: requires mathematical and source review.
\documentclass[11pt]{article}
\usepackage[margin=25mm]{geometry}
\usepackage{fontspec}
\setmainfont{Latin Modern Roman}
\usepackage{amsmath,amssymb,mathtools,unicode-math}
\setmathfont{Latin Modern Math}
\usepackage{enumitem,needspace,xurl,hyperref}
\setlength{\parindent}{0pt}
\setlength{\parskip}{5pt}
\setlength{\emergencystretch}{3em}
\newcommand{\N}{\mathbb{N}}
\newcommand{\Z}{\mathbb{Z}}
\newcommand{\Q}{\mathbb{Q}}
\newcommand{\R}{\mathbb{R}}
\newcommand{\C}{\mathbb{C}}
\newcommand{\F}{\mathbb{F}}
\newcommand{\A}{\mathbb{A}}
\newcommand{\PP}{\mathbb{P}}
\newcommand{\E}{\mathbb{E}}
\newcommand{\eps}{\varepsilon}
\newcommand{\dd}{\,\mathrm{d}}
\begin{document}
\section*{UnsolvedMath ID: 30002268}
\textbf{Source problem code:} catalogue entry
\subsection{Definitions and mathematical statement}
For a finite, nonempty, simple undirected graph \(G=(V,E)\), define its chromatic number by
\[
\chi(G)=\min\bigl\{k\in\N:k\geq 1,\ \exists c:V\to\{1,\ldots,k\}\ \forall\{u,v\}\in E,\ c(u)\neq c(v)\bigr\}.
\]
Define its maximum degree and clique number by
\[
\Delta(G)=\max_{v\in V}|\{u\in V:\{u,v\}\in E\}|,
\qquad
\omega(G)=\max\bigl\{|S|:S\subseteq V,\ \forall u\neq v\in S,\ \{u,v\}\in E\bigr\}.
\]
These definitions provide a provisional interpretation: the supplied record says only ``every graph'' and does not explicitly specify finiteness, simplicity, directedness, or conventions for the graph with no vertices.

Under this provisional interpretation, the literal supplied question is whether every such graph satisfies
\[
\chi(G)\leq\frac{\Delta(G)+1+\omega(G)}{2}.
\]
The record attributes this question to Reed's conjecture. Its displayed inequality contains no ceiling. The intended rounding requires source review: the alternative inequality
\[
\chi(G)\leq\left\lceil\frac{\Delta(G)+1+\omega(G)}{2}\right\rceil
\]
is a different assertion when the numerator is odd. No replacement of the supplied inequality is adopted here. The intended graph class and the attribution of the precise unrounded assertion remain unresolved.
\subsection{Short English statement}
The supplied question asks whether every graph has chromatic number at most half the sum of its maximum degree, its clique number, and one. The graph domain and the absence of a ceiling require review before this can be accepted as a precise statement of the attributed conjecture.
\subsection{Sources}
\begin{itemize}
\item Local problem registry: supplied complete record for permanent UnsolvedMath ID \(30002268\), source problem code \texttt{catalogue entry}, titled ``Reed's Chromatic Number Conjecture,'' in the Oberwolfach Reports Open Problems collection. The record attributes the conjecture to Reed but supplies no proposer entry or original publication citation. Its proposed-year field is \(2013\); the role of that year requires clarification.
\item Supplied report citation: \emph{Graph Theory} (2013), Oberwolfach Reports, DOI 10.4171/owr/2013/02, pp. 67--128. \url{https://doi.org/10.4171/owr/2013/02}. The background describes this collection as drawn from workshop reports and their open-problem discussions.
\item Supplied literature reference: \emph{Colouring Graphs with Sparse Neighbourhoods}. \url{https://doi.org/10.1016/j.jctb.2022.01.009}.
\item Supplied literature reference: \emph{Graph Theory, Oberwolfach Report 02/2013}. \url{https://doi.org/10.4171/owr/2013/02}.
\item Literature assessment dated exactly \texttt{2026-08-21}: the supplied curation labels the problem ``partially solved'' and states that Reed's exact chromatic bound remains open for general graphs, while fractional, epsilon, and several structural cases have been proved. This is supplied dated curation, not independently verified here; the record itself requests independent verification of time-sensitive status and cited claims.
\item The registry labels the statement \texttt{corrected\_verified} and reports that its mathematical content was checked against the official report, with notation, context, and question form normalized without changing the source-level claim. This is a supplied verification claim, not verification performed for this draft, and does not resolve the rounding issue.
\end{itemize}
\end{document}
